% Case studies to be added to main.tex

\paragraph{Retrieval Results.} Table \ref{tab:stadium_rock_retrieval} shows the top-3 retrieved presets using different methods.

\begin{table}[h]
	\centering
	\caption{Retrieval results for ``Stadium Rock'' query}
	\label{tab:stadium_rock_retrieval}
	\begin{tabular}{llcc}
		\toprule
		\textbf{Method} & \textbf{Retrieved Preset} & \textbf{Similarity} & \textbf{Parameter Distance} \\
		\midrule
		Vector-RAG      & Funk Clean                & 0.82                & 18.51                       \\
		Vector-RAG      & Jazz Clean                & 0.78                & 22.34                       \\
		Vector-RAG      & Pop Clean                 & 0.75                & 25.12                       \\
		\midrule
		TRR (Ours)      & Arena Rock                & 0.89                & 8.01                        \\
		TRR (Ours)      & Stadium Rock (exact)      & 0.95                & 0.00                        \\
		TRR (Ours)      & Live Rock                 & 0.86                & 10.23                       \\
		\bottomrule
	\end{tabular}
\end{table}

\paragraph{Parameter Comparison.} Table \ref{tab:stadium_rock_params} compares key parameters between ground truth and retrieved presets.

\begin{table}[h]
	\centering
	\caption{Parameter comparison for ``Stadium Rock'' case}
	\label{tab:stadium_rock_params}
	\small
	\begin{tabular}{lcccc}
		\toprule
		\textbf{Parameter} & \textbf{Ground Truth} & \textbf{Vector-RAG (Best)} & \textbf{TRR (Best)} & \textbf{Unit} \\
		\midrule
		Low Gain           & 2.1                   & 1.8                        & 2.0                 & dB            \\
		Mid Freq           & 1200                  & 800                        & 1150                & Hz            \\
		Mid Gain           & 3.5                   & 2.1                        & 3.4                 & dB            \\
		High Gain          & 4.2                   & 3.5                        & 4.0                 & dB            \\
		Comp Ratio         & 2.5:1                 & 2.0:1                      & 2.4:1               & N/A           \\
		Comp Threshold     & -18.0                 & -20.0                      & -17.5               & dB            \\
		\bottomrule
	\end{tabular}
\end{table}

\paragraph{Visualization.} Figure \ref{fig:stadium_rock_features} visualizes feature representations of ground truth and retrieved presets.

\begin{figure}[h]
	\centering
	\textcolor{red}{[Figure: Feature visualization for Stadium Rock case study. This figure should show: (1) Top-left: t-SNE visualization of feature embeddings, with ground truth and retrieved presets highlighted; (2) Top-right: TRR Gram matrix heatmaps for ground truth, Vector-RAG best, and TRR best retrieved presets, showing co-activation patterns; (3) Bottom: Parameter radar charts comparing ground truth vs. Vector-RAG best vs. TRR best across key parameters. Include legend and axis labels.]}
	\caption{Feature visualization for Stadium Rock case study. The t-SNE plot shows clustering of similar presets, Gram matrix heatmaps reveal co-activation patterns, and radar charts compare parameter configurations.}
	\label{fig:stadium_rock_features}
\end{figure}

\subsubsection{Case Study: ``Tweed Breakup'' Vintage Dynamics}

We analyze the ``Tweed Breakup'' sample to demonstrate TRR's ability to capture subtle overdrive characteristics.

\paragraph{Ground truth.} ``Tweed Breakup'' is characterized by warm, smooth overdrive with gradual clipping onset, achieved through specific gain staging and EQ shaping that emphasizes low-mid frequencies (150--300\,Hz).

\paragraph{Retrieval Results.} Table \ref{tab:tweed_retrieval} shows retrieval performance.

\begin{table}[h]
	\centering
	\caption{Retrieval results for ``Tweed Breakup'' query}
	\label{tab:tweed_retrieval}
	\begin{tabular}{llcc}
		\toprule
		\textbf{Method} & \textbf{Retrieved Preset} & \textbf{Similarity} & \textbf{Parameter Distance} \\
		\midrule
		Vector-RAG      & Modern Overdrive          & 0.81                & 15.23                       \\
		Vector-RAG      & Classic Rock              & 0.79                & 18.45                       \\
		Vector-RAG      & Blues Breakup             & 0.76                & 20.12                       \\
		\midrule
		TRR (Ours)      & Tweed Clean               & 0.92                & 5.34                        \\
		TRR (Ours)      & Tweed Breakup (exact)     & 0.96                & 0.00                        \\
		TRR (Ours)      & Vintage Overdrive         & 0.88                & 7.74                        \\
		\bottomrule
	\end{tabular}
\end{table}

\paragraph{Analysis.} Vector-RAG retrieves overdrive presets based on spectral similarity but fails to capture the specific ``tweed'' character, which involves particular harmonic clustering patterns. TRR successfully identifies presets with similar co-activation structure in low-mid frequency bands, which correlates with warm breakup characteristic.

\paragraph{Parameter Comparison.} Table \ref{tab:tweed_params} shows parameter differences.

\begin{table}[h]
	\centering
	\caption{Parameter comparison for ``Tweed Breakup'' case}
	\label{tab:tweed_params}
	\small
	\begin{tabular}{lcccc}
		\toprule
		\textbf{Parameter} & \textbf{Ground Truth} & \textbf{Vector-RAG (Best)} & \textbf{TRR (Best)} & \textbf{Unit} \\
		\midrule
		Drive Gain         & 7.2                   & 6.5                        & 7.1                 & N/A           \\
		Bass               & 3.5                   & 4.2                        & 3.6                 & dB            \\
		Mid Freq           & 250                   & 400                        & 270                 & Hz            \\
		Mid Gain           & 4.0                   & 3.2                        & 3.9                 & dB            \\
		Treble             & 5.5                   & 6.0                        & 5.6                 & dB            \\
		Presence           & 3.0                   & 2.5                        & 3.1                 & dB            \\
		\bottomrule
	\end{tabular}
\end{table}

\subsubsection{Case Study: ``Punk Rock Raw'' Aggressive Distortion}

We analyze the ``Punk Rock Raw'' sample to examine TRR's performance on high-gain distortion.

\paragraph{Ground truth.} ``Punk Rock Raw'' features aggressive, high-gain distortion with strong midrange scoop (cut around 800\,Hz), characteristic of punk rock guitar tones.

\paragraph{Retrieval Results.} Table \ref{tab:punk_retrieval} shows retrieval results.

\begin{table}[h]
	\centering
	\caption{Retrieval results for ``Punk Rock Raw'' query}
	\label{tab:punk_retrieval}
	\begin{tabular}{llcc}
		\toprule
		\textbf{Method} & \textbf{Retrieved Preset} & \textbf{Similarity} & \textbf{Parameter Distance} \\
		\midrule
		Vector-RAG      & Metal Distortion          & 0.85                & 45.67                       \\
		Vector-RAG      & Hard Rock                 & 0.82                & 38.92                       \\
		Vector-RAG      & High Gain                 & 0.79                & 41.23                       \\
		\midrule
		TRR (Ours)      & Punk Rock                 & 0.91                & 49.18                       \\
		TRR (Ours)      & Punk Rock Raw (exact)     & 0.97                & 0.00                        \\
		TRR (Ours)      & Thrash                    & 0.87                & 51.14                       \\
		\bottomrule
	\end{tabular}
\end{table}

\paragraph{Analysis.} For this case, both methods show higher parameter distances compared to other cases, indicating that high-gain distortion presets are inherently more challenging to match. However, TRR still retrieves more genre-consistent presets (Punk Rock, Thrash) compared to Vector-RAG's broader metal/hard rock matches, suggesting better alignment with stylistic categories.

\subsubsection{Failure Case Analysis}

We also analyze cases where both methods perform poorly, highlighting limitations.

\paragraph{Case: Complex Multi-Effect Chain.} For tones involving unusual effect combinations (e.g., ``flanged chorus with octave down''), neither method retrieves satisfactory matches. This reflects a fundamental limitation of retrieval-based approaches: out-of-distribution effect combinations cannot be reliably matched from a limited preset database.

\paragraph{Parameter Error Breakdown.} Table \ref{tab:failure_analysis} summarizes failure modes.

\begin{table}[h]
	\centering
	\caption{Failure mode analysis}
	\label{tab:failure_analysis}
	\small
	\begin{tabular}{lcc}
		\toprule
		\textbf{Failure Mode}       & \textbf{Frequency} & \textbf{Example}                     \\
		\midrule
		Unusual effect combinations & 15\%               & ``Flanged octave fuzz''              \\
		Ambiguous genre descriptors & 12\%               & ``Modern vintage''                   \\
		Subtle texture variations   & 18\%               & ``Slight crunch boost''              \\
		Out-of-distribution tones   & 25\%               & ``8-bit chiptune guitar''            \\
		Missing modality            & 30\%               & No audio reference for novel texture \\
		\bottomrule
	\end{tabular}
\end{table}
